\SourceLeaf{157}
Fährt man auf dieselbe Art weiter fort, so erhält $\Omega$ die Gestalt
\[
\Omega = (\nu C)^2+(\nu' C')^2+(\nu'' C'')^2+\cdots+(\nu^{(n-1)} C^{(n-1)})^2+(\beta^{(n)} B^{(n)})^2,
\]
wo alle $C$ von $x$ frei sind.



\Needspace{6\baselineskip}
\begin{center}
\textsc{Gauss} an \textsc{Olbers}. Göttingen, 14. April 1819.
\end{center}

\begin{center}
. . . . . . . . . . . . . . . .
\end{center}

Eben so ist mir Ihre Erinnerung gegen die richtige Formel für das Gewicht
des Resultats einer Combination von 2 Bestimmungen, deren Gewichte $p$, $q$
sind, nemlich $\dfrac{pq}{p+q}$, nicht recht erklärlich, es sei denn, dass
Sie meinen Brief nicht mit der Stelle bei \textsc{Littrow} zusammengehalten haben.
Allerdings setze ich voraus, dass die partiellen Resultate von einander
\textls{unabhängig} sind. Aber es kommt darauf an, von \emph{was für} Combinationen
die Rede ist, und da ist in dem \emph{vorliegenden Fall} der Umstand, dass
$\dfrac{pq}{p+q}$ immer $<p$ und $<q$ wird, nicht ungereimt, sondern das
Gegentheil würde ungereimt sein. Nehmen wir einen andern ganz ähnlichen
Fall. \textsc{Schumacher} hat die Absicht einen Längengradbogen zu messen. Es seien
seine äussersten Punkte $A$ und $C$. Er kann dieselben, d. h. seine
Raketen, nicht unmittelbar vergleichen, sondern bedarf eines
Zwischenpunkts $B$; vergleicht er zuerst $A$ mit $B$ 100mal, und nachher
$C$ mit $B$ auch 100mal, so haben die partiellen Längenunterschiede
$B-A$, $C-B$ eine grosse, und gleich grosse Zuverlässigkeit, deren Gewicht
durch $100=p$, $100=q$ angezeigt wird. Allein die aus der Verbindung von
beiden abgeleitete Differenz der Örter $C$, $A$ wird ja offenbar weniger
zuverlässig sein, als wenn sie aus 100 unmittelbaren Vergleichungen
abgeleitet wäre. Wie viel das Gewicht des Resultats aus dieser Combination
geringer sei, kann nur die Wahrscheinlichkeitstheorie lehren. Sie ergibt,
dass dieses \textls{geschlossene} Resultat nur so genau sei, als das Mittel
aus 50 directen Vergleichungen. Nur wenn $q=\infty$, also $C-B$ absolut
genau, wird das \textls{geschlossene} Resultat für $C-A$ eben so genau, wie
das benutzte $B-A$ war, wie es ja auch die Natur der Sache mit sich
bringt. \textsc{Littrows} falsche Formel würde für $q=\infty$ das absurde Resultat
$4p$ geben. --- Ich vermuthe, dass \textsc{Littrow} dadurch irre geleitet ist, dass der
